% =========================================================================
%  Experiments section for the NFG-SS paper.
%  Drop-in TeX: % =========================================================================
%  Experiments section for the NFG-SS paper.
%  Drop-in TeX: % =========================================================================
%  Experiments section for the NFG-SS paper.
%  Drop-in TeX: % =========================================================================
%  Experiments section for the NFG-SS paper.
%  Drop-in TeX: \input{experiments.tex} from your main file.
%  Placeholders are marked \todo{...} — fill them in before submission.
% =========================================================================

\section{Experiments}
\label{sec:experiments}

We empirically validate \texttt{NFG-SS} on a non-convex distributed image
classification task. The goal of this section is twofold: (i) verify
that removing the full-gradient anchor does not hurt convergence in
practice, and (ii) measure the gain from exploiting second-order
similarity at a fixed communication budget.

\subsection{Setting}
\label{sec:exp-setup}

\paragraph{Task and architecture.} We train ResNet-18 (adapted to
$32{\times}32$ inputs) with cross-entropy loss on CIFAR-10. The
training set is partitioned uniformly at random across $n=11$ nodes:
one node is treated as the server (it holds $f_1$) and the remaining
$10$ act as clients $f_2,\dots,f_{11}$. We use the standard CIFAR-10
train/test split and hold out $5\%$ of the training set as a
validation split used only for hyperparameter selection. All methods
share the same data partition, the same initialization, and the same
evaluation protocol.

\paragraph{Communication metric.} Throughout this section we report
test accuracy as a function of the number of $\mathbb{R}^d$ vectors
transmitted between the server and the clients (sum of both
directions), which is the cost measure used in our analysis
(Section~\ref{sec:related-work}).

\subsection{Methods}
\label{sec:exp-methods}

We compare \texttt{NFG-SS} (Algorithm~\ref{alg:nfgss}) against
\texttt{SVRS}~\citep{}\todo{cite [22]} and \texttt{FedAvg}~\citep{}\todo{cite [24]}.
\texttt{SVRS} is run in an \emph{identical} setting to \texttt{NFG-SS}
— same model, partition, inexact proximal solver, and number of outer
epochs — so that the two algorithms differ \emph{only} in the
construction of the gradient estimator $v_t^{(s)}$, isolating the
effect of removing the full-gradient anchor. We honour the original
\texttt{SVRS} schedule literally: at the start of every outer round
\texttt{SVRS} performs an exact full-gradient anchor refresh
$g_{\mathrm{ref}}=\tfrac{1}{n}\sum_i\nabla(f_i-f_1)(w_{\mathrm{ref}})$
by querying \emph{every} client over its entire local dataset. The
corresponding $\mathcal{O}(n)$ vectors per refresh are counted toward
the communication budget on the horizontal axis of
Figure~\ref{fig:cifar10-curves}.

\texttt{FedAvg}, which has no notion of similarity-aware variance
reduction, is run in a \emph{communication-matched} regime: the total
number of transmitted vectors over training is set equal to that of
\texttt{NFG-SS}. This is deliberately a generous baseline (it relaxes
\texttt{FedAvg}'s per-round full-aggregation schedule) and serves to
confirm that any gain we observe is driven by the similarity structure
rather than by partial-participation savings.

The proximal step on $f_1$ in both \texttt{NFG-SS} and \texttt{SVRS}
is approximated by the same inexact solver — a fixed number $K$ of
stochastic-gradient iterations on the proximal objective with gradient
clipping; details in Appendix~\ref{appendix:prox-solver}.

\subsection{Hyperparameters}
\label{sec:exp-hyper}

For every method we tune the hyperparameters listed in
Table~\ref{tab:hyper-grid} via a TPE sampler with median pruning;
$N_{\mathrm{trials}}=\todo{N}$ trials per method, $\todo{T}$ warmup
epochs, pruning interval $1$ epoch. The score of a trial is its best
validation accuracy. We then evaluate the best configuration on the
held-out test set. Selected values are listed in
Appendix~\ref{appendix:hyperparams}.

\begin{table}[t]
\centering
\caption{Hyperparameter search grids.}
\label{tab:hyper-grid}
\todo{[авторы — заполнить грид]}
\begin{tabular}{lll}
\toprule
Method & Hyperparameter & Search values \\
\midrule
\multirow{4}{*}{\texttt{NFG-SS}}
  & step size $\theta$           & \todo{...} \\
  & client batch size $b$        & \todo{...} \\
  & inner-prox iterations $K$    & \todo{...} \\
  & inner-prox learning rate     & \todo{...} \\
\midrule
\multirow{3}{*}{\texttt{SVRS}}
  & step size $\theta$           & \todo{...} \\
  & client batch size $b$        & \todo{...} \\
  & inner-prox learning rate     & \todo{...} \\
\midrule
\multirow{2}{*}{\texttt{FedAvg}}
  & local learning rate          & \todo{...} \\
  & number of local steps        & \todo{...} \\
\bottomrule
\end{tabular}
\end{table}

\subsection{Results}
\label{sec:exp-results}

Figure~\ref{fig:cifar10-curves} reports test accuracy versus the
number of transmitted vectors. \texttt{NFG-SS} attains $\approx 0.81$
test accuracy, dominating \texttt{SVRS}, which plateaus around
$\approx 0.78$ in the same communication budget;
communication-matched \texttt{FedAvg} reaches $\todo{...}$. Beyond the
final accuracy, \texttt{NFG-SS} also reaches the \texttt{SVRS}
plateau substantially earlier in training, indicating that the
running-mean estimator $\tilde v_t^{(s)}$ provides variance control
comparable to a full-gradient anchor without the corresponding
$\mathcal{O}(n)$ synchronisation cost at the start of every epoch.
The gap to communication-matched \texttt{FedAvg} confirms that
exploiting second-order similarity is the dominant factor: at the
same communication budget a similarity-agnostic method, even one
given the ``unfair'' freedom to ignore its original full-aggregation
schedule, does not catch up.

\begin{figure}[t]
\centering
\todo{[прикрепить график test accuracy vs transmitted vectors]}
\caption{Test accuracy on CIFAR-10 as a function of the number of
transmitted client$\leftrightarrow$server vectors. \texttt{NFG-SS}
matches or exceeds the final accuracy of \texttt{SVRS} while never
invoking a full-participation round; communication-matched
\texttt{FedAvg} trails both similarity-based methods. The
$\mathcal{O}(n)$ vectors of every \texttt{SVRS} anchor refresh are
included in its horizontal-axis position. Curves are averaged over
\todo{$R$} random seeds; shaded regions denote min/max range.}
\label{fig:cifar10-curves}
\end{figure}

\paragraph{Takeaways.}
(i) Replacing the full-gradient anchor of \texttt{SVRS} with the
running-mean estimator $\tilde v$ from
\eqref{eq:tildev-recursion} \emph{improves} both final accuracy and
communication efficiency on CIFAR-10, in line with our
$\mathcal{O}(n\delta^2/\varepsilon^2)$ communication bound.
(ii) The gain over a generous, communication-matched \texttt{FedAvg}
is large and consistent across the training horizon, showing that
under second-order similarity it is possible to recover nearly the
full benefit of variance reduction without ever paying the
$\mathcal{O}(n)$ per-anchor synchronisation cost.

\subsection{Ablation: visiting only $K$ clients per epoch}
\label{sec:exp-ablation}

The convergence analysis covers the case of disjoint batches that
together cover all $n-1$ clients within one epoch. We additionally
study a more aggressive regime in which only $K<n-1$ clients are
visited per epoch, drawn as a non-overlapping slice of a persistent
random permutation of $\{2,\ldots,n\}$ that is reshuffled once
exhausted; the carry-over estimator is rescaled by $(n-1)/K$ to
compensate for the missing mass. This ablation tests robustness of
\texttt{NFG-SS} to even tighter per-epoch client budgets; full results
are deferred to Appendix~\ref{appendix:clip-ablation}.
 from your main file.
%  Placeholders are marked \todo{...} — fill them in before submission.
% =========================================================================

\section{Experiments}
\label{sec:experiments}

We empirically validate \texttt{NFG-SS} on a non-convex distributed image
classification task. The goal of this section is twofold: (i) verify
that removing the full-gradient anchor does not hurt convergence in
practice, and (ii) measure the gain from exploiting second-order
similarity at a fixed communication budget.

\subsection{Setting}
\label{sec:exp-setup}

\paragraph{Task and architecture.} We train ResNet-18 (adapted to
$32{\times}32$ inputs) with cross-entropy loss on CIFAR-10. The
training set is partitioned uniformly at random across $n=11$ nodes:
one node is treated as the server (it holds $f_1$) and the remaining
$10$ act as clients $f_2,\dots,f_{11}$. We use the standard CIFAR-10
train/test split and hold out $5\%$ of the training set as a
validation split used only for hyperparameter selection. All methods
share the same data partition, the same initialization, and the same
evaluation protocol.

\paragraph{Communication metric.} Throughout this section we report
test accuracy as a function of the number of $\mathbb{R}^d$ vectors
transmitted between the server and the clients (sum of both
directions), which is the cost measure used in our analysis
(Section~\ref{sec:related-work}).

\subsection{Methods}
\label{sec:exp-methods}

We compare \texttt{NFG-SS} (Algorithm~\ref{alg:nfgss}) against
\texttt{SVRS}~\citep{}\todo{cite [22]} and \texttt{FedAvg}~\citep{}\todo{cite [24]}.
\texttt{SVRS} is run in an \emph{identical} setting to \texttt{NFG-SS}
— same model, partition, inexact proximal solver, and number of outer
epochs — so that the two algorithms differ \emph{only} in the
construction of the gradient estimator $v_t^{(s)}$, isolating the
effect of removing the full-gradient anchor. We honour the original
\texttt{SVRS} schedule literally: at the start of every outer round
\texttt{SVRS} performs an exact full-gradient anchor refresh
$g_{\mathrm{ref}}=\tfrac{1}{n}\sum_i\nabla(f_i-f_1)(w_{\mathrm{ref}})$
by querying \emph{every} client over its entire local dataset. The
corresponding $\mathcal{O}(n)$ vectors per refresh are counted toward
the communication budget on the horizontal axis of
Figure~\ref{fig:cifar10-curves}.

\texttt{FedAvg}, which has no notion of similarity-aware variance
reduction, is run in a \emph{communication-matched} regime: the total
number of transmitted vectors over training is set equal to that of
\texttt{NFG-SS}. This is deliberately a generous baseline (it relaxes
\texttt{FedAvg}'s per-round full-aggregation schedule) and serves to
confirm that any gain we observe is driven by the similarity structure
rather than by partial-participation savings.

The proximal step on $f_1$ in both \texttt{NFG-SS} and \texttt{SVRS}
is approximated by the same inexact solver — a fixed number $K$ of
stochastic-gradient iterations on the proximal objective with gradient
clipping; details in Appendix~\ref{appendix:prox-solver}.

\subsection{Hyperparameters}
\label{sec:exp-hyper}

For every method we tune the hyperparameters listed in
Table~\ref{tab:hyper-grid} via a TPE sampler with median pruning;
$N_{\mathrm{trials}}=\todo{N}$ trials per method, $\todo{T}$ warmup
epochs, pruning interval $1$ epoch. The score of a trial is its best
validation accuracy. We then evaluate the best configuration on the
held-out test set. Selected values are listed in
Appendix~\ref{appendix:hyperparams}.

\begin{table}[t]
\centering
\caption{Hyperparameter search grids.}
\label{tab:hyper-grid}
\todo{[авторы — заполнить грид]}
\begin{tabular}{lll}
\toprule
Method & Hyperparameter & Search values \\
\midrule
\multirow{4}{*}{\texttt{NFG-SS}}
  & step size $\theta$           & \todo{...} \\
  & client batch size $b$        & \todo{...} \\
  & inner-prox iterations $K$    & \todo{...} \\
  & inner-prox learning rate     & \todo{...} \\
\midrule
\multirow{3}{*}{\texttt{SVRS}}
  & step size $\theta$           & \todo{...} \\
  & client batch size $b$        & \todo{...} \\
  & inner-prox learning rate     & \todo{...} \\
\midrule
\multirow{2}{*}{\texttt{FedAvg}}
  & local learning rate          & \todo{...} \\
  & number of local steps        & \todo{...} \\
\bottomrule
\end{tabular}
\end{table}

\subsection{Results}
\label{sec:exp-results}

Figure~\ref{fig:cifar10-curves} reports test accuracy versus the
number of transmitted vectors. \texttt{NFG-SS} attains $\approx 0.81$
test accuracy, dominating \texttt{SVRS}, which plateaus around
$\approx 0.78$ in the same communication budget;
communication-matched \texttt{FedAvg} reaches $\todo{...}$. Beyond the
final accuracy, \texttt{NFG-SS} also reaches the \texttt{SVRS}
plateau substantially earlier in training, indicating that the
running-mean estimator $\tilde v_t^{(s)}$ provides variance control
comparable to a full-gradient anchor without the corresponding
$\mathcal{O}(n)$ synchronisation cost at the start of every epoch.
The gap to communication-matched \texttt{FedAvg} confirms that
exploiting second-order similarity is the dominant factor: at the
same communication budget a similarity-agnostic method, even one
given the ``unfair'' freedom to ignore its original full-aggregation
schedule, does not catch up.

\begin{figure}[t]
\centering
\todo{[прикрепить график test accuracy vs transmitted vectors]}
\caption{Test accuracy on CIFAR-10 as a function of the number of
transmitted client$\leftrightarrow$server vectors. \texttt{NFG-SS}
matches or exceeds the final accuracy of \texttt{SVRS} while never
invoking a full-participation round; communication-matched
\texttt{FedAvg} trails both similarity-based methods. The
$\mathcal{O}(n)$ vectors of every \texttt{SVRS} anchor refresh are
included in its horizontal-axis position. Curves are averaged over
\todo{$R$} random seeds; shaded regions denote min/max range.}
\label{fig:cifar10-curves}
\end{figure}

\paragraph{Takeaways.}
(i) Replacing the full-gradient anchor of \texttt{SVRS} with the
running-mean estimator $\tilde v$ from
\eqref{eq:tildev-recursion} \emph{improves} both final accuracy and
communication efficiency on CIFAR-10, in line with our
$\mathcal{O}(n\delta^2/\varepsilon^2)$ communication bound.
(ii) The gain over a generous, communication-matched \texttt{FedAvg}
is large and consistent across the training horizon, showing that
under second-order similarity it is possible to recover nearly the
full benefit of variance reduction without ever paying the
$\mathcal{O}(n)$ per-anchor synchronisation cost.

\subsection{Ablation: visiting only $K$ clients per epoch}
\label{sec:exp-ablation}

The convergence analysis covers the case of disjoint batches that
together cover all $n-1$ clients within one epoch. We additionally
study a more aggressive regime in which only $K<n-1$ clients are
visited per epoch, drawn as a non-overlapping slice of a persistent
random permutation of $\{2,\ldots,n\}$ that is reshuffled once
exhausted; the carry-over estimator is rescaled by $(n-1)/K$ to
compensate for the missing mass. This ablation tests robustness of
\texttt{NFG-SS} to even tighter per-epoch client budgets; full results
are deferred to Appendix~\ref{appendix:clip-ablation}.
 from your main file.
%  Placeholders are marked \todo{...} — fill them in before submission.
% =========================================================================

\section{Experiments}
\label{sec:experiments}

We empirically validate \texttt{NFG-SS} on a non-convex distributed image
classification task. The goal of this section is twofold: (i) verify
that removing the full-gradient anchor does not hurt convergence in
practice, and (ii) measure the gain from exploiting second-order
similarity at a fixed communication budget.

\subsection{Setting}
\label{sec:exp-setup}

\paragraph{Task and architecture.} We train ResNet-18 (adapted to
$32{\times}32$ inputs) with cross-entropy loss on CIFAR-10. The
training set is partitioned uniformly at random across $n=11$ nodes:
one node is treated as the server (it holds $f_1$) and the remaining
$10$ act as clients $f_2,\dots,f_{11}$. We use the standard CIFAR-10
train/test split and hold out $5\%$ of the training set as a
validation split used only for hyperparameter selection. All methods
share the same data partition, the same initialization, and the same
evaluation protocol.

\paragraph{Communication metric.} Throughout this section we report
test accuracy as a function of the number of $\mathbb{R}^d$ vectors
transmitted between the server and the clients (sum of both
directions), which is the cost measure used in our analysis
(Section~\ref{sec:related-work}).

\subsection{Methods}
\label{sec:exp-methods}

We compare \texttt{NFG-SS} (Algorithm~\ref{alg:nfgss}) against
\texttt{SVRS}~\citep{}\todo{cite [22]} and \texttt{FedAvg}~\citep{}\todo{cite [24]}.
\texttt{SVRS} is run in an \emph{identical} setting to \texttt{NFG-SS}
— same model, partition, inexact proximal solver, and number of outer
epochs — so that the two algorithms differ \emph{only} in the
construction of the gradient estimator $v_t^{(s)}$, isolating the
effect of removing the full-gradient anchor. We honour the original
\texttt{SVRS} schedule literally: at the start of every outer round
\texttt{SVRS} performs an exact full-gradient anchor refresh
$g_{\mathrm{ref}}=\tfrac{1}{n}\sum_i\nabla(f_i-f_1)(w_{\mathrm{ref}})$
by querying \emph{every} client over its entire local dataset. The
corresponding $\mathcal{O}(n)$ vectors per refresh are counted toward
the communication budget on the horizontal axis of
Figure~\ref{fig:cifar10-curves}.

\texttt{FedAvg}, which has no notion of similarity-aware variance
reduction, is run in a \emph{communication-matched} regime: the total
number of transmitted vectors over training is set equal to that of
\texttt{NFG-SS}. This is deliberately a generous baseline (it relaxes
\texttt{FedAvg}'s per-round full-aggregation schedule) and serves to
confirm that any gain we observe is driven by the similarity structure
rather than by partial-participation savings.

The proximal step on $f_1$ in both \texttt{NFG-SS} and \texttt{SVRS}
is approximated by the same inexact solver — a fixed number $K$ of
stochastic-gradient iterations on the proximal objective with gradient
clipping; details in Appendix~\ref{appendix:prox-solver}.

\subsection{Hyperparameters}
\label{sec:exp-hyper}

For every method we tune the hyperparameters listed in
Table~\ref{tab:hyper-grid} via a TPE sampler with median pruning;
$N_{\mathrm{trials}}=\todo{N}$ trials per method, $\todo{T}$ warmup
epochs, pruning interval $1$ epoch. The score of a trial is its best
validation accuracy. We then evaluate the best configuration on the
held-out test set. Selected values are listed in
Appendix~\ref{appendix:hyperparams}.

\begin{table}[t]
\centering
\caption{Hyperparameter search grids.}
\label{tab:hyper-grid}
\todo{[авторы — заполнить грид]}
\begin{tabular}{lll}
\toprule
Method & Hyperparameter & Search values \\
\midrule
\multirow{4}{*}{\texttt{NFG-SS}}
  & step size $\theta$           & \todo{...} \\
  & client batch size $b$        & \todo{...} \\
  & inner-prox iterations $K$    & \todo{...} \\
  & inner-prox learning rate     & \todo{...} \\
\midrule
\multirow{3}{*}{\texttt{SVRS}}
  & step size $\theta$           & \todo{...} \\
  & client batch size $b$        & \todo{...} \\
  & inner-prox learning rate     & \todo{...} \\
\midrule
\multirow{2}{*}{\texttt{FedAvg}}
  & local learning rate          & \todo{...} \\
  & number of local steps        & \todo{...} \\
\bottomrule
\end{tabular}
\end{table}

\subsection{Results}
\label{sec:exp-results}

Figure~\ref{fig:cifar10-curves} reports test accuracy versus the
number of transmitted vectors. \texttt{NFG-SS} attains $\approx 0.81$
test accuracy, dominating \texttt{SVRS}, which plateaus around
$\approx 0.78$ in the same communication budget;
communication-matched \texttt{FedAvg} reaches $\todo{...}$. Beyond the
final accuracy, \texttt{NFG-SS} also reaches the \texttt{SVRS}
plateau substantially earlier in training, indicating that the
running-mean estimator $\tilde v_t^{(s)}$ provides variance control
comparable to a full-gradient anchor without the corresponding
$\mathcal{O}(n)$ synchronisation cost at the start of every epoch.
The gap to communication-matched \texttt{FedAvg} confirms that
exploiting second-order similarity is the dominant factor: at the
same communication budget a similarity-agnostic method, even one
given the ``unfair'' freedom to ignore its original full-aggregation
schedule, does not catch up.

\begin{figure}[t]
\centering
\todo{[прикрепить график test accuracy vs transmitted vectors]}
\caption{Test accuracy on CIFAR-10 as a function of the number of
transmitted client$\leftrightarrow$server vectors. \texttt{NFG-SS}
matches or exceeds the final accuracy of \texttt{SVRS} while never
invoking a full-participation round; communication-matched
\texttt{FedAvg} trails both similarity-based methods. The
$\mathcal{O}(n)$ vectors of every \texttt{SVRS} anchor refresh are
included in its horizontal-axis position. Curves are averaged over
\todo{$R$} random seeds; shaded regions denote min/max range.}
\label{fig:cifar10-curves}
\end{figure}

\paragraph{Takeaways.}
(i) Replacing the full-gradient anchor of \texttt{SVRS} with the
running-mean estimator $\tilde v$ from
\eqref{eq:tildev-recursion} \emph{improves} both final accuracy and
communication efficiency on CIFAR-10, in line with our
$\mathcal{O}(n\delta^2/\varepsilon^2)$ communication bound.
(ii) The gain over a generous, communication-matched \texttt{FedAvg}
is large and consistent across the training horizon, showing that
under second-order similarity it is possible to recover nearly the
full benefit of variance reduction without ever paying the
$\mathcal{O}(n)$ per-anchor synchronisation cost.

\subsection{Ablation: visiting only $K$ clients per epoch}
\label{sec:exp-ablation}

The convergence analysis covers the case of disjoint batches that
together cover all $n-1$ clients within one epoch. We additionally
study a more aggressive regime in which only $K<n-1$ clients are
visited per epoch, drawn as a non-overlapping slice of a persistent
random permutation of $\{2,\ldots,n\}$ that is reshuffled once
exhausted; the carry-over estimator is rescaled by $(n-1)/K$ to
compensate for the missing mass. This ablation tests robustness of
\texttt{NFG-SS} to even tighter per-epoch client budgets; full results
are deferred to Appendix~\ref{appendix:clip-ablation}.
 from your main file.
%  Placeholders are marked \todo{...} — fill them in before submission.
% =========================================================================

\section{Experiments}
\label{sec:experiments}

We empirically validate \texttt{NFG-SS} on a non-convex distributed image
classification task. The goal of this section is twofold: (i) verify
that removing the full-gradient anchor does not hurt convergence in
practice, and (ii) measure the gain from exploiting second-order
similarity at a fixed communication budget.

\subsection{Setting}
\label{sec:exp-setup}

\paragraph{Task and architecture.} We train ResNet-18 (adapted to
$32{\times}32$ inputs) with cross-entropy loss on CIFAR-10. The
training set is partitioned uniformly at random across $n=11$ nodes:
one node is treated as the server (it holds $f_1$) and the remaining
$10$ act as clients $f_2,\dots,f_{11}$. We use the standard CIFAR-10
train/test split and hold out $5\%$ of the training set as a
validation split used only for hyperparameter selection. All methods
share the same data partition, the same initialization, and the same
evaluation protocol.

\paragraph{Communication metric.} Throughout this section we report
test accuracy as a function of the number of $\mathbb{R}^d$ vectors
transmitted between the server and the clients (sum of both
directions), which is the cost measure used in our analysis
(Section~\ref{sec:related-work}).

\subsection{Methods}
\label{sec:exp-methods}

We compare \texttt{NFG-SS} (Algorithm~\ref{alg:nfgss}) against
\texttt{SVRS}~\citep{}\todo{cite [22]} and \texttt{FedAvg}~\citep{}\todo{cite [24]}.
\texttt{SVRS} is run in an \emph{identical} setting to \texttt{NFG-SS}
— same model, partition, inexact proximal solver, and number of outer
epochs — so that the two algorithms differ \emph{only} in the
construction of the gradient estimator $v_t^{(s)}$, isolating the
effect of removing the full-gradient anchor. We honour the original
\texttt{SVRS} schedule literally: at the start of every outer round
\texttt{SVRS} performs an exact full-gradient anchor refresh
$g_{\mathrm{ref}}=\tfrac{1}{n}\sum_i\nabla(f_i-f_1)(w_{\mathrm{ref}})$
by querying \emph{every} client over its entire local dataset. The
corresponding $\mathcal{O}(n)$ vectors per refresh are counted toward
the communication budget on the horizontal axis of
Figure~\ref{fig:cifar10-curves}.

\texttt{FedAvg}, which has no notion of similarity-aware variance
reduction, is run in a \emph{communication-matched} regime: the total
number of transmitted vectors over training is set equal to that of
\texttt{NFG-SS}. This is deliberately a generous baseline (it relaxes
\texttt{FedAvg}'s per-round full-aggregation schedule) and serves to
confirm that any gain we observe is driven by the similarity structure
rather than by partial-participation savings.

The proximal step on $f_1$ in both \texttt{NFG-SS} and \texttt{SVRS}
is approximated by the same inexact solver — a fixed number $K$ of
stochastic-gradient iterations on the proximal objective with gradient
clipping; details in Appendix~\ref{appendix:prox-solver}.

\subsection{Hyperparameters}
\label{sec:exp-hyper}

For every method we tune the hyperparameters listed in
Table~\ref{tab:hyper-grid} via a TPE sampler with median pruning;
$N_{\mathrm{trials}}=\todo{N}$ trials per method, $\todo{T}$ warmup
epochs, pruning interval $1$ epoch. The score of a trial is its best
validation accuracy. We then evaluate the best configuration on the
held-out test set. Selected values are listed in
Appendix~\ref{appendix:hyperparams}.

\begin{table}[t]
\centering
\caption{Hyperparameter search grids.}
\label{tab:hyper-grid}
\todo{[авторы — заполнить грид]}
\begin{tabular}{lll}
\toprule
Method & Hyperparameter & Search values \\
\midrule
\multirow{4}{*}{\texttt{NFG-SS}}
  & step size $\theta$           & \todo{...} \\
  & client batch size $b$        & \todo{...} \\
  & inner-prox iterations $K$    & \todo{...} \\
  & inner-prox learning rate     & \todo{...} \\
\midrule
\multirow{3}{*}{\texttt{SVRS}}
  & step size $\theta$           & \todo{...} \\
  & client batch size $b$        & \todo{...} \\
  & inner-prox learning rate     & \todo{...} \\
\midrule
\multirow{2}{*}{\texttt{FedAvg}}
  & local learning rate          & \todo{...} \\
  & number of local steps        & \todo{...} \\
\bottomrule
\end{tabular}
\end{table}

\subsection{Results}
\label{sec:exp-results}

Figure~\ref{fig:cifar10-curves} reports test accuracy versus the
number of transmitted vectors. \texttt{NFG-SS} attains $\approx 0.81$
test accuracy, dominating \texttt{SVRS}, which plateaus around
$\approx 0.78$ in the same communication budget;
communication-matched \texttt{FedAvg} reaches $\todo{...}$. Beyond the
final accuracy, \texttt{NFG-SS} also reaches the \texttt{SVRS}
plateau substantially earlier in training, indicating that the
running-mean estimator $\tilde v_t^{(s)}$ provides variance control
comparable to a full-gradient anchor without the corresponding
$\mathcal{O}(n)$ synchronisation cost at the start of every epoch.
The gap to communication-matched \texttt{FedAvg} confirms that
exploiting second-order similarity is the dominant factor: at the
same communication budget a similarity-agnostic method, even one
given the ``unfair'' freedom to ignore its original full-aggregation
schedule, does not catch up.

\begin{figure}[t]
\centering
\todo{[прикрепить график test accuracy vs transmitted vectors]}
\caption{Test accuracy on CIFAR-10 as a function of the number of
transmitted client$\leftrightarrow$server vectors. \texttt{NFG-SS}
matches or exceeds the final accuracy of \texttt{SVRS} while never
invoking a full-participation round; communication-matched
\texttt{FedAvg} trails both similarity-based methods. The
$\mathcal{O}(n)$ vectors of every \texttt{SVRS} anchor refresh are
included in its horizontal-axis position. Curves are averaged over
\todo{$R$} random seeds; shaded regions denote min/max range.}
\label{fig:cifar10-curves}
\end{figure}

\paragraph{Takeaways.}
(i) Replacing the full-gradient anchor of \texttt{SVRS} with the
running-mean estimator $\tilde v$ from
\eqref{eq:tildev-recursion} \emph{improves} both final accuracy and
communication efficiency on CIFAR-10, in line with our
$\mathcal{O}(n\delta^2/\varepsilon^2)$ communication bound.
(ii) The gain over a generous, communication-matched \texttt{FedAvg}
is large and consistent across the training horizon, showing that
under second-order similarity it is possible to recover nearly the
full benefit of variance reduction without ever paying the
$\mathcal{O}(n)$ per-anchor synchronisation cost.

\subsection{Ablation: visiting only $K$ clients per epoch}
\label{sec:exp-ablation}

The convergence analysis covers the case of disjoint batches that
together cover all $n-1$ clients within one epoch. We additionally
study a more aggressive regime in which only $K<n-1$ clients are
visited per epoch, drawn as a non-overlapping slice of a persistent
random permutation of $\{2,\ldots,n\}$ that is reshuffled once
exhausted; the carry-over estimator is rescaled by $(n-1)/K$ to
compensate for the missing mass. This ablation tests robustness of
\texttt{NFG-SS} to even tighter per-epoch client budgets; full results
are deferred to Appendix~\ref{appendix:clip-ablation}.
